\documentclass[11pt,a4paper]{article}
\usepackage[T1]{fontenc}
\usepackage[margin=2.5cm]{geometry}
\usepackage{xcolor}
\usepackage[hidelinks]{hyperref}
\newcommand{\pending}[1]{\textcolor{red}{\textbf{[PENDING: #1]}}}
\pagestyle{empty}
\begin{document}
\noindent Anandharajan TRV\\
Department of Electronics and Communication Engineering\\
St.\ Joseph College of Engineering\\
Sriperumbudur, Chennai 602117, India\\[1em]
\noindent \today\\[1.5em]

\noindent The Editors\\
\emph{Computational Biology and Chemistry}\\[1.5em]

\noindent Dear Editors,\\[0.8em]

I am submitting ``Local--edge routing for conformational B-cell epitope
prediction: what survives leakage-free selection'' for consideration as a
research article.

The paper does two things. It characterises, over ten seeds with confidence
intervals, what a laboratory gains by routing antigens between a local
predictor and a remote one according to the local model's uncertainty: roughly
a third of the latency of always-remote execution and a large reduction in how
many antigen sequences leave the device, with no quality difference resolved
either way at the fifteen antigens available. The same
accounting shows that routing does \emph{not} reduce modelled carbon --- it is
the most carbon-intensive policy in our cost model, because the model that decides
whether to route is not cheap relative to the one it gates. That contradicts
the usual framing of edge inference as an environmental measure, and the
mechanism is general enough to matter beyond this application.

The second half is a negative result, and I want to be direct with you about
it. An earlier, unpublished analysis in this project found that fusing our
local predictor with BepiPred-3.0 improved on BepiPred-3.0 alone. In preparing
this manuscript I found that the fusion weight and the decision threshold had
both been selected on the same split the result was reported on. The paper
reproduces those earlier figures under their original protocol, to within the
precision at which the original grid was stored, then re-estimates the same
quantity with both parameters chosen by nested leave-one-sequence-out
selection. The advantage is no longer detectable, and the MCC estimate changes
sign. The datasets are too small to show that fusion cannot help, and I do not
claim that; what they do show is how far the selection step moves the
estimate. I report this rather than quietly dropping it, because the selection
step involved is easy to commit, hard to see in a write-up, and sufficient here
to turn a null result into a plausible-looking positive one.

The work fits the journal's scope in computational life sciences: its subject
is how conformational B-cell epitope predictors are selected and evaluated,
and its primary contribution is a leakage-free re-estimation of a published
style of fusion result. It follows directly from Attique et al.,
``DeepBCE'', published in this journal in 2023, which motivated the project.
I make no claim to state-of-the-art predictive accuracy and say so explicitly;
our standalone predictor remains below BepiPred-3.0 on BP3C50ID.

All code, the raw per-seed results, and a single script that regenerates every
number, table and figure in the manuscript are available at
\url{https://github.com/Anandharajan/edgeconbce}, and will be archived on
Zenodo with a DOI at acceptance.

I am the sole author and the corresponding author. The manuscript is original,
is not under consideration elsewhere, and I have approved it for submission.
Funding, competing interests and the use of generative AI tools are declared in
the manuscript.

\vspace{1em}
\noindent Yours sincerely,\\[1.5em]
\noindent Anandharajan TRV

\end{document}
